\section{The higher-rank classification}
\textbf{Theorem S.Classification.} Let S=S(t,k,v) be a finite Steiner
system with 3≤t\textless k, and let G=Aut(S). Then Γ(S) is a firm
residually connected geometry on which G is flag-transitive if and only
if S is one of the following:

\begin{enumerate}
\def\labelenumi{\arabic{enumi}.}
\tightlist
\item
  the points and affine planes of AG(3,2);
\item
  the projective GF(4)-subline design on PG(1,16);
\item
  one of the three large Witt designs S(3,6,22), S(4,7,23), S(5,8,24).
\end{enumerate}

Each resulting geometry also satisfies (IP)₂. In the first case the full
design group is AGL(3,2); in the second it is PΓL(2,16). In the last
three cases the full design groups contain the indicated native Mathieu
groups, and their orders are 887,040,10,200,960,244,823,040 respectively.

\textbf{Proof.} Suppose first that Γ(S) is firm and G is
flag-transitive. The firmness criterion gives k−t+1≥2 and
r=(v−t+1)/(k−t+1)≥3. The equivariant chamber bijection of Theorem S.3 makes
G t-transitive on the points of S. Kantor's Theorem3, retained with its
precise primary-source hypotheses, therefore leaves the binary affine
plane family, the projective subline family, and the five Mathieu
designs.

For the binary affine plane family \(S_d\), the one-point derivative is
PG(d−1,2). The residue theorem makes Γ(\(S_d\)) residually connected
exactly when its rank-three derived construction is connected.
B.ProjectivePlane includes the positive case PG(2,2), while
B.SpaceObstruction excludes every projective dimension at least three.
Thus d=3. The case d=2 does not satisfy the already necessary condition
r≥3. Conversely the explicit affine action in the binary affine argument
above is chamber-transitive, and the same residue argument proves
residual connectivity when d=3. Preservation of the fourth point of each
affine plane proves Aut(\(S_d\))=AGL(d,2), as proved above.

In the projective subline case the parameters are S(3,q+1,qᵉ+1), with
e≥2 and q≥3. The derivative at a point is AG(e,q). B.AffinePlane and
B.SpaceObstruction, together with the residue theorem, imply residual
connectivity exactly when e=2. For this remaining case Theorem I.Subline
proves Aut(S)=PΓL(2,q²). If q=pᶠ, the chamber count is
\textbar PGL(2,q²)\textbar q whereas the full group order is
\textbar PGL(2,q²)\textbar·2f. Flag transitivity therefore requires q to
divide 2f.

This divisibility leaves only q=4 in the range q≥3. If f=1, it is
impossible. If f≥3, then q≥2ᶠ\textgreater2f, also impossible. If f=2,
q=p² divides 4 only when p=2. At q=4 the normalization in that theorem
proves transitivity explicitly: the relevant Frobenius action cycles
through the four nontrivial direction classes and the scalar action
covers the three eligible points in each class. This supplies
sufficiency as well as necessity. Firmness and (IP)₂ follow from the
uniform Steiner theorems.

For the two small Witt designs the chamber counts from Theorem S.3
are 23,760 and285,120. The literal finite models and full automorphism
computations give orders 7,920 and95,040. A transitive group cannot have
an orbit larger than itself, so neither design gives a flag-transitive
construction.

For each large Witt design r=5. The explicit finite maps identify every
(t−2)-fold derivative with PG(2,4): one derivative is verified against
coordinates, and native t-transitivity transports it to every other
derivative. B.ProjectivePlane and the residue theorem give residual
connectivity. The parameters give firmness, and Theorem S.6 gives (IP)₂.
The already retained Mathieu fiber computation verifies transitivity of
a native subgroup on every stage of the chamber fiber; consequently the
full automorphism group is also transitive. This proves sufficiency and
completes the classification. ∎

The source's theorem already states this final list. The contribution
here is a deduction that separates the common structural proofs,
retained classification inputs, and exact finite fibers, and that
continues to prove the correct final list after correcting the source's
intermediate residual-connectivity assertions. It is not a new
classification theorem.
